\documentclass[../SQG_notes.tex]{subfiles}

\begin{document}
    \section{Literature Review}
    
    \subsection{Ryan et.al.}
    \textbf{Year: 2025} \par
       The $\QGone$ model incorporates the first-order corrections that were neglected in the basic QG approximation. It essentially refines the QG equations by accounting for non-geostrophic (ageostrophic) flow components that are dependent on the Rossby number ($\epsilon$). \par
    In \textbf{Chapter 2}, the $\text{QG}^{+1}$ model is introduced. In this paper, 
    \[ N = f \equiv \text{Constant} \mn{page 8}\]
    To facilitate the asymptotic approximation, a potential field is introduced. 
    \[ \mathbf{A} = (-G, -F, \Phi)\]
    By Incompressible condition we have 
    \[ \mathbf{v} = \nabla_3 \times \mathbf{A} \mn{In this paper $\nabla_3$ is 3D gradient. 2D is just $\nabla$}\]
    Some Physical implications of the model
    \begin{enumerate}
        \item Breaking Symmetry of QG model.
        \item It Captures Cyclogeostrophic balance. 
        \begin{center}
        \scriptsize
        \noindent
        Cyclogeostrophic balance is a fundamental force balance approximation used in meteorology and physical oceanography to describe the motion of fluids (like air and water) in curved paths, where the Coriolis force is balanced by the pressure gradient force and the centrifugal force. It is an essential extension of the simpler geostrophic balance, which only considers straight flow. This balance is particularly important in systems with high curvature and strong winds, such as tropical cyclones (hurricanes/typhoons), mid-latitude low-pressure systems, and strong ocean eddies. The Governing Equation is 
        \begin{equation}
            \underbrace{f \mathbf{v}}_{\text{Coriolis Force}} + \underbrace{\frac{|\mathbf{v}|^2}{R}}_{\text{Centrifugal Force}} = \underbrace{-\frac{1}{\rho} \frac{\partial p}{\partial n}}_{\text{Pressure Gradient Force}}
        \end{equation}
        Here $n$ is the normal direction pointing toward the center of curvature.
        \end{center}
        \item Inclusiong of \textbf{Frontogenesis} \mn{Generation of Ocean Fronts}. 
    \end{enumerate}
    In \textbf{Chapter 3.} A simulation for $\QGone$ is conducted, showing several features:
    \begin{enumerate}
        \item More Vigorous due to captureing ageostrophic frontogenesis.
        \item Since the Ageostrophic effects creates stronger surface velocity. Finer structure can be seen on surface using $\QGone$. \mn{See Figure 4 in page 24}.
    \end{enumerate}
    In summary, this paper provides a very detailed derivation to the $\QGone$ equation which is introduced more detailed in \ref{sec:QGone model}. This paper also demonstrate two simulation to show how $\QGone$ model captures balanced submesoscale dynamics and frontogenesis.

    \begin{remark}
        In my own derivation following Ryan's work in the later chapter, first half of the calcualtion is based on this paper, then I follow Ryan's note for the rest. 
    \end{remark}

    \subsubsection{Questions Regrading this Paper}

    \begin{enumerate}
        \item In page 10, Equation 15. How does the above derivation implies $F = G = 0$. Does this implies that 
        \[ F^0 + \epsilon F^1 + \cdots = 0\]
        Then each order term of $F$ and $G$ is 0. I just don't get the physical meaning here. 
    \end{enumerate}

    \subsection{J.Wang et.al. Reconstructing the Ocean's Interior from Surface Data}
    \textbf{Year : 2013} \par

    In the \textbf{Introduction}, the author discussed the current challenge of using SSH and SST \mn{Surface Sea Height and Surface Sea Temperature} measurement to reconstruct subsurface dynamics. \par
    \begin{itemize}
        \item     Traditional studies assume the signal is dominated by barotropic and first baroclinic modes. However, these modes are typically calculated by \textbf{assuming buoyancy anomalies vanish at the surface}. 
        \item SQG theorys works as well. But it normaly assume $0$ interior PV. 
    \end{itemize}
    The author introduced the \textbf{Interior plus surface QG} method. It is quasigeostrophic. As introduced in Chapter 2: 
    \begin{enumerate}
        \item Surface buoyancy anomoly contributes to the surface part of streamfunction $\psi^s$ : 
        \[
        \begin{gathered}
        \mathcal{L} \Psi+f_0+\beta y=Q, \\
        \mathcal{L}=\left(\frac{\partial^2}{\partial x^2}+\frac{\partial^2}{\partial y^2}+\frac{\partial}{\partial z} \frac{f_0^2}{N^2} \frac{\partial}{\partial z}\right), \quad \text { and }-H<z<0,
        \end{gathered}
        \]
        Governing equation is 
        \[
        \mathcal{L} \psi^s=0
        \]
        Essentially, this is same in the SQG theory where we assume $0$ interior PV. With boundary condition : 
        \[
        \frac{\partial}{\partial z} \psi^s(\mathbf{x}, z, t)=b(\mathbf{x}, z, t) / f_0 \quad \text { at } \quad z=0,-H,
        \]
        \item The interior part is governed by 
        \[
        \frac{\partial}{\partial z} \frac{f_0^2}{N^2} \frac{\partial}{\partial z} \hat{\psi}^i-\kappa^2 \hat{\psi}^i=\hat{q}^i \quad \text { with } \quad \frac{d \hat{\psi}^i}{d z}=0 \quad \text { at } \quad z=0,-H .
        \]
    \end{enumerate}
    However, the interior $q^i$ is not clear. 
    \begin{remark}
        This is the essential modification of this isQG model. They project the interior induced PV equation onto baroclinic modes and \textbf{impose additional boundary conditions to deduce the gravest modes.}
    \end{remark}

    % \subsection{D.J. Muraki, C. Snyder & R. Rotunno 1999}

    % This is the origin of $\QGone$ theory. 

    \subsection{John R. Taylor and Andrew F. Thompson 2023}

    This review focuses on submesosclae dynamics in the upper ocean with a particular focus on multiscale interactions involving the submesoscale. 

    \subsubsection{Generation of Submesoscale Flow}
    Submesosclaes are energized by a avariety of instabilities that develop from balanced currents. 

    \subsubsection{Energy Transfer and Multiscale Interactions}
    \subsubsection*{Down-Scale Connections}
    Energy enters the ocean at large scales through surface wind. The mesocslae baroclinic instability then converts the energy to low vertical modes. Submesosclae motions are potential mechanisms for supplying energy from ocean mososcale to small-scale turbulence and scales that ultimately dissipated. \mn{this indicates that submesoscale motion is one of the driving mechanisms in energy dissipation in ocean.}

    Simulation shows that the formation of fronts is one of such path. 

    \subsection{Ross et al. 2023}
    To coarse-grain a high-resolution field $\hat{q}_\kappa$ onto a low-resolution grid, Ross et al. propose three filtering and coarse-graining operators. Let $\kappa$ denote the wavenumber, $\Delta x_{\text{lores}}$ the grid spacing of the coarse model, and define the cutoff $\kappa^c \equiv 2/3$ of the low-resolution Nyquist frequency.

    \subsubsection{Operator 1: Spectral Truncation with Sharp Filter}
    The first option applies the same quadruple-exponential filter used by \texttt{pyqg} for small-scale dissipation. It leaves small wavenumbers unchanged and attenuates wavenumbers above $\kappa^c$:
    \begin{equation}
        \hat{\bar{q}}_\kappa =
        \begin{cases}
            \hat{q}_\kappa, & \kappa < \kappa^c, \\
            \hat{q}_\kappa \cdot e^{-23.6(\kappa - \kappa^c)^4 \Delta x_{\text{lores}}^4}, & \kappa \geq \kappa^c.
        \end{cases}
    \end{equation}
    This is the most conservative choice, closest to no filtering at all, since it is already applied within the ocean model. \mn{``Operator 1'' = spectral truncation followed by this filter.}

    \subsubsection{Operator 2: Spectral Truncation with Softer Gaussian Filter}
    The second option applies a Gaussian filter to all remaining modes:
    \begin{equation}
        \hat{\bar{q}}_\kappa = \hat{q}_\kappa \cdot e^{-\kappa^2 (2 \Delta x_{\text{lores}})^2 / 24}.
    \end{equation}
    This choice follows Guan et al. (2022) and Pope (2000). By the filter-width definition of Lund (1997), this filter is twice as large as the grid size of the coarse model.

    \subsubsection{Operator 3: Diffusion-Based Filtering with Real-Space Coarsening}
    The third option is closer to the procedure used in ocean models. It applies \texttt{GCM-Filters} (Grooms et al. 2021; Loose et al. 2022), which approximates the spectral transfer function of a Gaussian filter using polynomials based on the Laplacian diffusion operator. The \texttt{pyqg} output is converted onto a Cartesian grid, and the filtered field is then coarse-grained in real space: to reduce the resolution by a factor $K$, the input field is averaged over non-overlapping $K \times K$ boxes.

    \begin{remark}
        A comparison of the effects of these three filtering and coarse-graining operators on PV and its subgrid forcing is shown in Figures 3 and 4 of Ross et al. (2023).
    \end{remark}

    \subsubsection{Performance Metrics (Chapter 4)}
    Since the number of candidate parameterizations is too large to inspect manually, Ross et al. define a hierarchy of metrics. They split into \textbf{offline} metrics (evaluated on held-out testing sets of subgrid forcing data) and \textbf{online} metrics (computed by actually running parameterized low-resolution simulations and comparing them against high-resolution truth).

    \paragraph{Offline Metrics.} Let $S$ denote the true subgrid forcing and $\hat{S}$ the parameterization's prediction. For each fluid layer:
    \begin{enumerate}
        \item \textbf{Coefficient of determination} ($R^2$):
        \begin{equation}
            R^2 = 1 - \frac{\mathbb{E}\!\left[(S - \hat{S})^2\right]}{\mathbb{E}\!\left[(S - \mathbb{E}[S])^2\right]}.
        \end{equation}
        $R^2 = 1$ for perfect predictions, $R^2 = 0$ when predictions are no better than the mean, and $R^2 < 0$ when worse than predicting the mean.

        \item \textbf{Pearson correlation} ($\rho$):
        \begin{equation}
            \rho = \frac{\mathrm{Cov}(S, \hat{S})}{\sigma_S \, \sigma_{\hat{S}}}.
        \end{equation}
        $\rho \in [-1, 1]$. \mn{$\rho$ can stay high even when $R^2$ is negative --- e.g., predictions off by a large but consistent scaling factor.}
    \end{enumerate}
    These are evaluated on filtered/coarse-grained high-resolution data from both eddy and jet configurations, aggregated over time and space (or expressed as functions thereof). Power and energy-redistribution spectra of the predicted forcing are also visualized against ground truth.

    \paragraph{Online Metrics.} A new low-resolution QG simulation is run for $10$ years (5 random seeds, both eddy and jet configurations) with the parameterization adding to the PV tendency at every time step. Distance metrics between low-resolution and high-resolution distributions are then computed and \emph{normalized by the corresponding distance for the unparameterized low-resolution simulation}, giving an interpretable similarity score.

    \begin{enumerate}
        \item \textbf{Spectral differences} (\texttt{spectral\_diff}). RMS difference between time-averaged power spectra and spectral fluxes:
        \begin{equation}
            \texttt{spectral\_diff}(\text{sim}_a, \text{sim}_b; f) \equiv \sqrt{\frac{1}{|\mathcal{K}|} \sum_{k \in \mathcal{K}} \left(f_{\text{sim}_a}(k) - f_{\text{sim}_b}(k)\right)^2},
        \end{equation}
        where $\mathcal{K}$ is a set of isotropic wavenumbers, log-spaced up to $2/3$ of the low-resolution Nyquist frequency. Computed for energy and enstrophy spectra in each layer and for spectral energy fluxes (KE flux, APE flux, APE generation, bottom drag) --- $8$ metrics total.

        \item \textbf{Distribution differences} (\texttt{distrib\_diff}). Earth mover's (Wasserstein) distance between the empirical PDFs of state variables at the end of the simulation:
        \begin{equation}
            \texttt{distrib\_diff}(\text{sim}_a, \text{sim}_b; f) \equiv \int_{-\infty}^{\infty} \left| P_{\text{sim}_a}(f \leq x) - P_{\text{sim}_b}(f \leq x) \right| \, dx.
        \end{equation}
        Computed for the marginal distributions of $u$, $v$, $q$, kinetic-energy density $(u^2 + v^2)/2$, and enstrophy $\frac{1}{2}(\nabla \times \mathbf{u})^2$ in each layer --- $10$ metrics total. \mn{Visualizes ``climate'' --- long-term statistics --- not ``weather.''}

        \item \textbf{Decorrelation time difference} (\texttt{decorr\_time}). The minimum time at which two simulations starting from nearby initial conditions decorrelate below a threshold $\delta$:
        \begin{equation}
            \texttt{decorr\_time}(\text{sim}_a, \text{sim}_b) \equiv \mathbb{E}_{q_0, \epsilon}\!\left[ \min_{t} \left\{ t : \mathrm{Corr}\!\left( q^{(t)}_{\text{sim}_a}(q_0), \, q^{(t)}_{\text{sim}_b}(q_0 + \epsilon) \right) \leq \delta \right\} \right],
        \end{equation}
        with $\delta = 0.5$ and $\epsilon$ a Gaussian perturbation of standard deviation $10^{-10}$. When the two simulations differ in resolution, the higher-resolution field is first filtered and coarse-grained using Operator 1 before the correlation is computed. \mn{Captures forecast skill --- ``weather'' --- rather than long-term climate.}
    \end{enumerate}

    \begin{remark}
        The three online metrics correspond to the three panels of Figure 5 in Ross et al.\ (2023): (a) \texttt{spectral\_diff} as RMSE in spectral space, (b) \texttt{distrib\_diff} as the area between two PDFs, (c) \texttt{decorr\_time} as how quickly a parameterized low-resolution run diverges from its high-resolution reference. Together they probe model physics, climatology, and short-term forecast skill, respectively.
    \end{remark}




\end{document}